\documentclass[11pt,a4paper]{article}
\usepackage[T1]{fontenc}
\usepackage[utf8]{inputenc}
\usepackage[french]{babel}
\usepackage{amsmath,amssymb,amsthm,mathtools}
\usepackage[margin=2.5cm]{geometry}
\usepackage[hidelinks]{hyperref}
\newtheorem{lemme}{Lemme T}
\newtheorem{theoreme}{Théorème de synthèse}
\newtheorem*{corollaire}{Corollaire expérimental certifié}
\newcommand{\ellg}{\ell}
\newcommand{\vten}{\nu_{10}}
\title{Syracuse $\Phi$--$\Psi$\\Lemmes dérivés sur les créations de tête et la superposition}
\author{Complément de travail aux théorèmes 5.19 et 5.37}
\date{17 novembre 2026}
\begin{document}\maketitle
\begin{abstract}
Cette note déduit des identités exactes sur les variations de longueur décimale,
les frontières de création temporaire ou persistante, la composition temporelle,
l'érosion des suffixes superposés et les jonctions après divergence de parité.
Les résultats sont universels à horizon fini. Le record initial
$N=274133054632352106267$ sert d'exemple de vérification, non de preuve
asymptotique ni de preuve de la conjecture de Syracuse.
\end{abstract}

\section*{Notations et base déjà établie}
La transformation classique et la transformation condensée sont respectivement
\[
C(n)=\begin{cases}n/2&n\equiv0\pmod2,\\3n+1&n\equiv1\pmod2,\end{cases}
\qquad
U(n)=\begin{cases}n/2&n\equiv0\pmod2,\\(3n+1)/2&n\equiv1\pmod2.\end{cases}
\]
Pour $n\ge1$, sa longueur décimale est $\ellg(n)=1+\lfloor\log_{10}n\rfloor$.
On prolonge $U$ à $0$ par $U(0)=0$.\par
Les cellules $\Phi$--$\Psi$ réalisent exactement $U$ sur tout état décimal
(théorème 5.19 de la v9). Le lemme 5.37 établit notamment la coïncidence
initiale des parités de $n$ et $n+a10^L$ jusqu'au rang
$r=L+\nu_2(a)$, où se produit la première divergence. Nous ne redémontrons
pas ces deux résultats, mais les utilisons explicitement.

\begin{lemme}[Pas de longueur élémentaire]
Pour tout $n\ge1$,
\[
\ellg(C(n))-\ellg(n)\in\{-1,0,+1\}.
\]
Si $n$ est impair, la variation vaut $0$ ou $+1$; s'il est pair,
elle vaut $0$ ou $-1$.
\end{lemme}
\begin{proof}
Si $\ellg(n)=L$, alors $10^{L-1}\le n<10^L$.
Pour $n$ impair, $n<3n+1<10^{L+1}$; pour $n$ pair,
$10^{L-1}/2\le n/2<10^L/2$. Ces bornes interdisent toute autre variation.
\end{proof}

\begin{lemme}[Porte décimale et persistance immédiate]
Soit $n$ impair de longueur $L$. Une création de tête par $3n+1$
a lieu si et seulement si
\[
3n+1\ge10^L\quad\Longleftrightarrow\quad
n\ge\frac{10^L-1}{3}=\underbrace{33\cdots3}_{L\text{ chiffres}}.
\]
La tête nouvellement créée disparaît à l'opération suivante de division
par deux si et seulement si
\[
10^L\le3n+1<2\cdot10^L.
\]
Elle subsiste après cette division si et seulement si $3n+1\ge2\cdot10^L$.
\end{lemme}
\begin{proof}
$3n+1<3\cdot10^L+1<10^{L+1}$ pour $L\ge1$ et $n<10^L$;
le franchissement de $10^L$ caractérise donc la création d'une unique position.
La parité de $3n+1$ est paire, et sa division conserve la nouvelle position
si et seulement si $(3n+1)/2\ge10^L$.
\end{proof}

\begin{lemme}[Bilan télescopique des têtes]
Soit $x_t=C^t(x_0)$, et notons
\[
A(s,t)=\#\{j\in[s,t):\ellg(x_{j+1})-\ellg(x_j)=1\},\quad
B(s,t)=\#\{j\in[s,t):\ellg(x_{j+1})-\ellg(x_j)=-1\}.
\]
Alors, pour tous $0\le s<t$,
\[
\boxed{\ellg(x_t)-\ellg(x_s)=A(s,t)-B(s,t).}
\]
Si $t_i$ et $t_{i+1}$ sont deux instants successifs de première atteinte
d'une longueur supérieure à toutes les longueurs précédentes,
alors $A(t_i,t_{i+1})-B(t_i,t_{i+1})=1$.
\end{lemme}
\begin{proof}
Il suffit de sommer $\ellg(x_{j+1})-\ellg(x_j)$ sur
$j=s,\ldots,t-1$. Le lemme T1 assure qu'aucune longueur record
ne peut être sautée.
\end{proof}

\begin{lemme}[Composition multiplicative exacte sur un intervalle]
Pour $x_j=C^j(x_0)$, $0\le s<t$, notons $p$ le nombre de pas impairs
$3x_j+1$, $q$ le nombre de divisions $x_j/2$, et
$\mathcal I=\{j\in[s,t):x_j\text{ impair}\}$.
Alors
\[
\boxed{\frac{x_t}{x_s}=\frac{3^p}{2^q}
\prod_{j\in\mathcal I}\left(1+\frac{1}{3x_j}\right).}
\]
En conséquence,
\[
\ellg(x_t)-\ellg(x_s)
=\left\lfloor\log_{10}x_s+p\log_{10}3-q\log_{10}2
+\sum_{j\in\mathcal I}\log_{10}\left(1+\frac1{3x_j}\right)\right\rfloor
-\lfloor\log_{10}x_s\rfloor.
\]
\end{lemme}
\begin{proof}
À chaque pas impair, $x_{j+1}/x_j=3(1+1/(3x_j))$;
à chaque pas pair, $x_{j+1}/x_j=1/2$.
Le produit des ratios successifs se télescope.
\end{proof}

\begin{lemme}[Érosion exacte du suffixe dans une superposition]
Soit $n\ge0$, $L\ge1$, $a\in\{1,\ldots,9\}$,
$n'=n+a10^L$ et $r=L+\nu_2(a)$.
Posons $x_h=U^h(n)$, $y_h=U^h(n')$ et
$m_h=\#\{0\le j<h:x_j\text{ impair}\}$.
Pour $0\le h\le r$, on a
\[
\boxed{y_h-x_h=\frac{a10^L3^{m_h}}{2^h}},\qquad
\boxed{\nu_2(y_h-x_h)=r-h}.
\]
En notant $\nu_{10}(z)=\max\{k\ge0:10^k\mid z\}$ pour $z\neq0$,
on obtient de plus la formule exacte
\[
\boxed{\vten(y_h-x_h)=\min\bigl(L,r-h\bigr),\qquad 0\le h\le r.}
\]
Le nombre de chiffres décimaux de pied identiques commence donc à $L$,
reste égal à $L$ pendant $\nu_2(a)$ transitions, puis décroît exactement
d'une unité par transition jusqu'à $0$ au rang $r$.
\end{lemme}
\begin{proof}
Comme $\Delta_0=a10^L$ est divisible par $2^r$ et pas par $2^{r+1}$,
la différence reste paire jusqu'au rang $r-1$. Les deux trajectoires suivent
donc les mêmes branches avant le rang $r$ et la première identité résulte
de la composition affine (argument du lemme 5.37).
En prenant les valuations, $\nu_2(y_h-x_h)=L+\nu_2(a)-h$ et
$\nu_5(y_h-x_h)=L+\nu_5(a)$.
Par définition, $\nu_{10}=\min(\nu_2,\nu_5)$.
Comme $a$ est un chiffre non nul, il ne peut être simultanément
multiple de $2$ et de $5$; cela réduit le minimum à $\min(L,r-h)$.
\end{proof}

\begin{lemme}[Franchissement de tête pendant la superposition]
Sous les hypothèses du lemme T5, pour toute position $k\ge1$ et
tout $0\le h\le r$, le nombre $y_h$ franchit la frontière décimale
$10^k$ alors que $x_h$ ne la franchit pas si et seulement si
\[
\boxed{x_h<10^k\le y_h
\ \Longleftrightarrow\ 
x_h<10^k\text{ et }a10^L3^{m_h}\ge2^h(10^k-x_h).}
\]
Cette condition peut être satisfaite bien avant la première divergence
de parité $h=r$.
\end{lemme}
\begin{proof}
Pour $h\le r$, substituer au terme $y_h-x_h$ l'identité exacte
du lemme T5. La comparaison avec $10^k$ est alors équivalente à
l'inégalité affichée, sans approximation.
\end{proof}

\begin{lemme}[Jonction de superposition après rupture de parité]
Pour deux trajectoires quelconques $x_h=U^h(n)$,
$y_h=U^h(n')$, posons $\Delta_h=y_h-x_h$ et
$p_h=x_h\bmod2$, $p'_h=y_h\bmod2$.
Alors, à tous les rangs $h\ge0$,
\[
\boxed{\Delta_{h+1}
=\frac{(1+2p'_h)\Delta_h+(p'_h-p_h)(2x_h+1)}2.}
\]
En particulier, lorsque $p_h=p'_h$, la différence suit la loi
multiplicative $\Delta_{h+1}=(1+2p_h)\Delta_h/2$;
à la première rupture, un terme de jonction dépendant de l'état $x_h$
apparaît exactement.
\end{lemme}
\begin{proof}
On utilise $U(z)=((1+2(z\bmod2))z+(z\bmod2))/2$ pour $z=x_h,y_h$,
puis on soustrait et on remplace $y_h=x_h+\Delta_h$.
\end{proof}

\begin{theoreme}[Couplage tête--pied et extension décimale à horizon fini]
Soit $N\ge1$. Pour chaque $L\ge1$, on définit
$S_L=N\bmod10^L$ et $a_L=\lfloor N/10^L\rfloor\bmod10$.
Si $a_L=0$, $S_{L+1}=S_L$ et leurs trajectoires sont identiques.
Si $a_L\neq0$, on pose $r_L=L+\nu_2(a_L)$.
Alors, pour tout horizon fini $H$, les différences
$U^h(S_{L+1})-U^h(S_L)$, leurs suffixes communs et les franchissements
de frontières $10^k$ sont entièrement déterminés par :
(1) le lemme T5 pour $h\le r_L$;
(2) le lemme T7 à partir de $h=r_L$;
(3) le critère de frontière T6 aux rangs communs;
(4) le même opérateur $\mathscr F$ du théorème 5.19 pour tous les états.
Aucune règle locale nouvelle n'est requise pour la variation de longueur
ou la jonction après rupture de parité.
\end{theoreme}
\begin{proof}
On procède par récurrence sur $h$. Avant et jusqu'à $r_L$, le lemme T5
fournit la différence et T6 permet de décider tout franchissement.
À partir de $r_L$, les deux états complets sont déterminés par
$\mathscr F$ et la différence au pas suivant par T7.
La récurrence fournit les identités à chaque horizon fini.
\end{proof}

\section*{Contrôle arithmétique du record à 42 chiffres}
Pour $N=274133054632352106267$, les 570 premières transitions de $C$
comprennent 247 pas impairs et 323 pas pairs : 114 montées locales de
longueur, 93 diminutions et 363 transitions sans variation de longueur.
D'où $114-93=21$, et $\ellg(C^{570}(N))=42$.
Les 21 premières atteintes des longueurs $22,\ldots,42$ sont toutes suivies
d'une division faisant perdre immédiatement le chiffre créé.
L'état au rang 570 est
\[
C^{570}(N)=113298124744388651798242538014293435290632.
\]
Pour la superposition des suffixes de 20 et 21 chiffres, $a_{20}=2$ et
$r_{20}=20+\nu_2(2)=21$.
Le lemme T5 donne donc exactement
\[
\nu_{10}\bigl(U^h(S_{21})-U^h(S_{20})\bigr)
=\min(20,21-h),\quad 0\le h\le21.
\]
Au rang classique $3$, la tête atteint 22 chiffres pour $N$ mais seulement
21 chiffres pour $S_{20}$; cette divergence des longueurs n'exige pas
une divergence préalable des parités.

\section*{Portée et limite logique}
T1--T7 constituent des identités et critères universels;
le contrôle sur $N$ est un exemple fini.
Le théorème de couplage complète les observables de la superposition
sans prétendre classifier les orbites à horizon infini.
Les préfixes finis de profondeurs dyadiques restent tous réalisables
(selon l'étude v15); aucune exclusion universelle des trajectoires
divergentes n'est démontrée ici.

\end{document}
